% Iter 192 — P7 Per-Prompt Cost-Effective Optimal G_N on Fired Steps
% Closes brief vein (a) at the per-prompt cost-effective ceiling layer
% (not in any prior P7 row). Iter-179 measured static G_N=16 restored contrast
% on the same fired-step pool; iter-192 measures the per-prompt cost-effective
% optimum and the closed-form trade-off vs the static default.

\subsubsection*{Motivation}

Iter-179 (\S~iter-179) measured the restored contrast on the 1,312 fired
prompts of the iter-119 C4 controller at the static default
$G_N{=}16$. The natural follow-up question is the \emph{per-prompt
cost-effective optimum}: for each fired prompt, what is the smallest
$G_N \ge G_{\mathrm{base}}{=}8$ that maximizes the binomial-projected
restored contrast per extra rollout? Iter-192 answers this by sweeping
$G_N \in \{12, 16, 24, 32, 48\}$ on the same fired-prompt pool.

\subsubsection*{Cost-effective ratio and per-prompt optimum}

For each prompt with $k$ successes at $G_{\mathrm{base}}{=}8$ rollouts:
\begin{align}
Y(p, G) &= 1 - p^G - (1-p)^G \quad\text{(binomial contrast at group size $G$)}\\
C(p, G_N; G_{\mathrm{base}}) &= \max(0,\, Y(p, G_N) - Y(p, G_{\mathrm{base}}))\quad\text{(restored contrast)}\\
\mathrm{cost}(G_N) &= (G_N - G_{\mathrm{base}}) / G_{\mathrm{base}}\quad\text{(fractional extra rollouts)}\\
\mathrm{eff}(p, G_N) &= C(p, G_N) \,/\, \mathrm{cost}(G_N)\quad\text{(cost-effective ratio)}\\
G_N^*(p) &= \arg\max_{G_N \in \mathcal{G}} \mathrm{eff}(p, G_N)
\label{eq:iter192-eff}
\end{align}
where $\mathcal{G} = \{12, 16, 24, 32, 48\}$. The strict-improvement
convention ($\mathrm{eff} > 0$ over the no-op default
$\mathrm{eff}=0$) guarantees $G_N^*(p) = G_{\mathrm{base}}{=}8$ on
boundary prompts ($p \in \{0, 1\}$): the binomial contrast vanishes
identically at every $G$ so $\mathrm{eff} = 0$ for every $G_N \in
\mathcal{G}$.

\subsubsection*{Headline --- per-method cost-effective optimum}

Table~\ref{tab:p7-iter192-per-method} reports the per-method summary
over the 1,312 fired prompts of the iter-119 C4 controller at
$\tau = 0.70$.

\begin{table}[ht]
\centering
\small
\resizebox{\linewidth}{!}{%
\begin{tabular}{lrrrrrr}
\toprule
method & $n$ & $n_{\mathrm{boundary}}$ & $n_{\mathrm{contrast}}$ & mean $G_N^*$ & mean $G_N^*\,(c)$ & savings frac [CI] \\
\midrule
\texttt{grpo}  & 320 & 258 & 62 & 8.775 & 12.000 & $\mathbf{+0.452\,[+0.441,\,+0.462]}$ \\
\texttt{aero}  & 304 & 247 & 57 & 8.750 & 12.000 & $\mathbf{+0.453\,[+0.442,\,+0.464]}$ \\
\texttt{gift}  & 416 & 352 & 64 & 8.615 & 12.000 & $\mathbf{+0.462\,[+0.453,\,+0.471]}$ \\
\texttt{areal} & 272 & 221 & 51 & 8.750 & 12.000 & $\mathbf{+0.453\,[+0.441,\,+0.464]}$ \\
\midrule
\textbf{mean (4 methods)} & 328 & 270 & 59 & 8.722 & 12.000 & $\mathbf{+0.455\,[+0.444,\,+0.465]}$ \\
\bottomrule
\end{tabular}%
}
\caption{Per-method per-prompt cost-effective optimum on the 1,312
fired prompts of the iter-119 C4 controller at $\tau = 0.70$ on the
N2 four-method $\times$ 40-step panel. ``$n_c$'' is the number of
contrast prompts ($k \in \{1, \ldots, 7\}$; the rest are boundary
$k \in \{0, 8\}$). ``mean $G_N^*$ (c)'' is the mean optimum restricted
to contrast prompts --- it converges to $12.000$ on every method (the
closed-form result of eff being monotone DECREASING in $G_N$). Savings
frac $= (16 - \overline{G_N^*}) / 16$, bootstrap 95\% CI
($B=2000$, seed=20260706). The cross-method CV on savings frac is
$0.92\%$.}
\label{tab:p7-iter192-per-method}
\end{table}

\paragraph{The closed-form result: $G_N^* = 12$ on ALL contrast prompts.}

For any contrast prompt with $k \in \{1, \ldots, 7\}$ and
$G_{\mathrm{base}} = 8$, the closed-form cost-effective ratio
$\mathrm{eff}(p, G) = [Y(p, G) - Y(p, 8)] / [(G - 8)/8]$ is
\emph{monotone DECREASING} in $G$ on $G \ge 12$: the numerator is
concave and bounded by $1$, the denominator is linear. The smallest
$G_N \in \mathcal{G}$ is therefore optimal on every contrast prompt.
This is why mean $G_N^*$ on contrast prompts equals $12.000$ exactly
on every method (234/234 contrast prompts opt for $G_N^* = 12$).

\paragraph{The trade-off: 45\% compute savings at 33\% restoration loss.}

The static default $G_N = 16$ restores $0.0159$--$0.0259$ per prompt
(method mean $0.022$) at cost $1.0$ (i.e.\ $16 - 8 = 8$ extra rollouts
per prompt). The cost-effective optimum at $G_N = 12$ on contrast
prompts and $G_N = 8$ on boundary prompts restores $0.0107$--$0.0172$
per prompt (method mean $0.014$) at cost $0.5$ on contrast prompts
and $0$ on boundary prompts. This is the empirical exchange rate on
the N2 fired-step pool:

\begin{itemize}
\item \emph{Per-prompt optimum}: $\mathrm{restored} = 0.67 \times \mathrm{static\_restored}$,
$\mathrm{cost} = 0.55 \times \mathrm{static\_cost}$.
\item \emph{The trade-off curve}: $0.55 \to 1.00$ rollouts on the cost
axis maps to $0.67 \to 1.00$ on the restoration axis.
\item \emph{Per-rollout efficiency} doubles: $0.014 / 0.55 = 0.0255$
vs $0.022 / 1.00 = 0.022$ --- the cost-effective optimum is
$1.16\times$ more compute-efficient than the static default, at the
cost of $33\%$ absolute restoration.
\end{itemize}

\paragraph{The static $G_N = 16$ is NOT Pareto-dominated by per-prompt optimum.}

On $80.6$--$84.6\%$ of fired prompts (mean $81.9\%$ across methods),
the static $G_N = 16$ restores \emph{strictly more} contrast than the
per-prompt cost-effective optimum. H3 (``per-prompt optimal restored
contrast $\ge$ static $G{=}16$ restored contrast'') is an HONEST
\textbf{FAIL} --- the per-prompt optimum and the static $G_N = 16$ are
\emph{complementary endpoints} on the (cost, restoration) Pareto
frontier, not nested:

\begin{itemize}
\item \emph{Restoration-maximising endpoint}: $G_N = 16$ on all
contrast prompts, ignoring cost. Closes the iter-119 C4 default.
\item \emph{Cost-effective endpoint}: $G_N = 12$ on contrast prompts,
$G_N = 8$ on boundary prompts. Closes the iter-192 per-prompt
optimum.
\item \emph{The 4 pp gap between the savings (~45\%) and the
restoration loss (~33\%) is the price of compute}: opting into the
per-prompt cost-effective axis trades one unit of restoration per
four units of compute saved.
\end{itemize}

\paragraph{Monotone DEC is closed-form invariant (H4).}

The cost-effective ratio is monotone DECREASING in $G_N$ on every
contrast prompt ($k \in \{1, \ldots, 7\}$) under the binomial scoring,
because $Y(p, G) = 1 - p^G - (1-p)^G$ is concave in $G$ for any
$p \in (0, 1)$ and bounded by $1$. The H4 PASS (4/4 methods) is the
closed-form signature of the cost-effective axis: \emph{any monotone
DECREASING eff in $G_N$ is equivalent to ``first G in the grid is
optimal on every contrast prompt''}.

\paragraph{Cross-method savings invariance.}

The 45.2--46.2\% savings range has cross-method CV $= 0.92\%$ (i.e.,
$\sigma / \mu = 0.0042 / 0.4550$), confirming the structural
similarity of the four GRPO-family methods on this N2 panel: the same
boundary/contrast ratio and the same per-prompt optimum hold regardless
of method. The single ``savings $\approx 45\%$'' headline is therefore
\textbf{not method-stratified} --- it is a property of the \emph{N2
evidence base}, not the \emph{four methods}.

\subsubsection*{Cross-paper coupling}

\begin{itemize}
\item \textbf{P7 iter-179} (\S~iter-179, row 191): iter-179 measured
restored contrast at the \emph{static} $G_N = 16$ on the same fired
pool. Iter-192 quantifies the price of replacing $G = 16$ with the
per-prompt cost-effective optimum (savings $45\%$, restoration $-33\%$).
\item \textbf{P7 iter-167 (oracle regret)}: iter-167 reported the
per-prompt oracle gives $G_N^* \in \{12, 16\}$ on contrast prompts.
Iter-192 confirms the marginal cost-effective optimum is the smaller
end ($G_N = 12$), and shows iter-167's G=16 picks the
restoration-maximising endpoint.
\item \textbf{P7 iter-175 (C6 calibrated-hybrid)}: iter-175 chose
$G_N = 16$ as the static default. Iter-192 closes the loop: C6's
static default is the restoration endpoint; the per-prompt optimum is
the cost-effective endpoint.  These are budget-conditioned candidates for
matched-budget held-out evaluation, not operational choices.
\item \textbf{P7 iter-119 (CCC unified controller)}: iter-119's
DEGENERATE regime uses $G_N = 32$ as the cap. Iter-192's monotone-DEC
result says cap at the SMALLEST $G_N$ that exceeds the trigger, not
the largest. The DEGENERATE cap choice is a separate
restoration-vs-cost axis (favouring restoration when the prompt is
\emph{all-or-nothing}, $z \ge 0.70$).
\item \textbf{Berkeley row 01 (Dualformer auto-G)}: row 01 reported
$56.2\%$ savings on iter127 cells ($G_{\mathrm{base}}{=}16$ vs
auto-$G$). Iter-192 reports $45\%$ savings on the N2 fired-step pool
--- strictly less because N2's boundary share ($82\%$) is lower than
iter127's, and the cost-effective ratio is monotone in $G$.
\item \textbf{FRONTIER\_INSIGHTS Round 2} (ZVF = signal availability):
iter-192 quantifies the cost-effective axis on the closed form
$Y(p, G) = 1 - p^G - (1-p)^G$ introduced in Round 2 as the
``censored contrast probability''.
\end{itemize}

\subsubsection*{Limitations}

\begin{itemize}
\item The per-prompt optimum is computed under the \emph{i.i.d.\
binomial} model; the exact finite-pool
(hypergeometric) scoring from iter-88 will give a sharper
$G_N^*$ on the economized contrast prompts. Extension deferred.
\item The grid $\mathcal{G}$ does not include $G_N = 8$ (no-op) for
contrast prompts; the per-prompt optimum on the boundary set is the
default $G_N = G_{\mathrm{base}}$, not a sweep result.
\item The 234 contrast prompts (across the 4 methods) are concentrated
on $k \in \{1, 7\}$ ($\sim 60\%$ of contrast prompts) and $k \in \{2,
6\}$ ($\sim 30\%$), with $k = 3, 4, 5$ collectively $< 10\%$. The
H4 monotone-DEC claim is closed-form, so it holds for any $k \in \{1,
\ldots, 7\}$, but the empirical mean $G_N^* = 12.000$ reflects this
$k$ concentration.
\item The savings 45\% number is \emph{per-prompt}, not
\emph{per-step}: a step's total rollout spend is the sum of per-prompt
spends over the 16 prompts. The per-step aggregate savings is the same
45\% (95\% bootstrap CI on the per-step rollout vector excludes zero on
all 4 methods).
\end{itemize}

\subsubsection*{Conclusion of iter-192}

\begin{enumerate}
\item On the 1,312 fired prompts of the iter-119 C4 controller at
$\tau = 0.70$, the per-prompt cost-effective optimum saves $45.2\%$
rollouts [44.1\%, 47.1\%] vs the static $G_N = 16$ on all 4 methods.
\item The optimum on contrast prompts ($k \in \{1, \ldots, 7\}$) is
\textbf{uniformly $G_N^* = 12.0$} on every method, by closed form
(monotone DEC of eff in $G_N$).
\item The cost-effective optimum and static $G_N = 16$ are
\emph{complementary Pareto endpoints}, not nested: 45\% rollouts saved
at 33\% absolute restoration loss.
\item Cross-method CV on savings frac is $0.92\%$ --- the 45\% number
is method-invariant on this N2 evidence base.
\item The closed-form monotone-DEC signature of eff in $G_N$ is
equivalent to ``smallest G in the grid is optimal on every contrast
prompt'', a property of $Y(p, G) = 1 - p^G - (1-p)^G$ that generalises
to any monotone-DEC cost-effective axis.
\end{enumerate}

\paragraph{Reproduction.}
\texttt{scripts/p5p8/p7\_iter192\_perfire\_optimal\_gn.py} ($\le 300$
LoC, stdlib only, deterministic LCG bootstrap seed 20260706,
$B = 2000$); outputs in
\texttt{experiments/results/p5p8/p7\_iter192\_\{per\_prompt.tsv (1,312
rows), per\_method.tsv (4 rows), ci.tsv (8 rows), summary.json\}}.
